% ==============================================================================
% CAPÍTULO 5 — DIRETRIZES DE ANONIMIZAÇÃO, PSEUDONIMIZAÇÃO E SANITIZAÇÃO
% ==============================================================================

\chapter{Diretrizes de Anonimização, Pseudonimização e Sanitização}
\label{cap:anonimizacao}

Este capítulo desenvolve as diretrizes técnicas para a desidentificação de dados pessoais em logs de infraestrutura, apresentando as técnicas reconhecidas pela ANPD e pelo GDPR e demonstrando sua aplicação prática por meio de uma ferramenta construída com o Microsoft Presidio, framework de código aberto para detecção e anonimização de PII \cite{microsoft2023presidio}, o \textit{MikroTik Log Anonymizer}, aplicada a logs reais de roteadores MikroTik RouterOS. A detecção combina três recognizers: o recognizer nativo do Presidio para endereços IP, um \textit{PatternRecognizer} customizado para endereços MAC -- inexistente no conjunto padrão do framework -- e padrões de expressão regular para nomes de usuário em eventos típicos de autenticação MikroTik.

% ─────────────────────────────────────────────────────────────────────────────
\section{Conceitos e Distinções}
\label{sec:conceitos-anon}

A literatura de privacidade e a legislação distinguem três operações de desidentificação com efeitos legais e técnicos distintos:

\begin{alineas}
    \item \textbf{Anonimização}: processo irreversível pelo qual o dado perde completamente a possibilidade de associação a um titular, mesmo com o uso de meios técnicos razoáveis. Dados devidamente anonimizados saem do escopo de aplicação da LGPD (art.\ 12);
    \item \textbf{Pseudonimização}: substituição de identificadores diretos por pseudônimos, mantendo a possibilidade de reversão mediante informação adicional mantida separadamente. O dado pseudonimizado continua sendo dado pessoal nos termos da LGPD, mas seu tratamento beneficia-se de salvaguardas reduzidas em determinados contextos (art.\ 13, §4°);
    \item \textbf{Sanitização}: remoção ou modificação de dados em um documento ou arquivo para reduzir riscos, sem necessariamente atingir o grau de anonimização completa. Inclui técnicas como supressão de campos, truncamento e mascaramento parcial.
\end{alineas}

% ─────────────────────────────────────────────────────────────────────────────
\section{Técnicas Recomendadas por Tipo de Dado}
\label{sec:tecnicas}

A escolha da técnica de anonimização deve considerar a finalidade do log após o tratamento: se será usado apenas para auditoria interna (onde correlação entre eventos é necessária) ou para compartilhamento externo (onde a identidade do titular é irrelevante). O \autoref{qdr:tecnicas} apresenta as técnicas recomendadas para cada tipo de dado pessoal presente em logs de infraestrutura.

\begin{quadro}[htb]
    \caption{\label{qdr:tecnicas}Técnicas de anonimização recomendadas por tipo de dado em logs}
    \begin{tabular}{|p{2.8cm}|p{3cm}|p{3.2cm}|p{4.8cm}|}
        \hline
        \textbf{Tipo de dado} & \textbf{Uso interno} & \textbf{Compartilhamento externo} & \textbf{Justificativa} \\
        \hline
        Endereço IP & Hash SHA-256 (pseudonimização) & Generalização ou supressão & Hash preserva correlação entre eventos sem expor o IP real \\
        \hline
        Endereço MAC & Mascaramento dos 3 primeiros octetos & Supressão completa & Os 3 primeiros octetos identificam o fabricante; os 3 últimos identificam o dispositivo \\
        \hline
        Nome de usuário & Pseudonimização por hash & Supressão & Necessário para correlação de ações do mesmo usuário em auditorias \\
        \hline
        Timestamp & Preservar (necessário) & Generalização (hora ou dia) & Precisão de segundo cria risco de quasi-identificação por combinação \\
        \hline
        Nome de host & Pseudonimização & Supressão & Pode revelar identidade organizacional ou pessoal \\
        \hline
        Porta TCP/UDP & Preservar & Preservar ou generalizar & Baixo risco isolado; relevante para análise de segurança \\
        \hline
    \end{tabular}
    \legend{Fonte: Própria dos autores (2025), com base em \citeonline{anpd2022anonimizacao}.}
\end{quadro}

% ─────────────────────────────────────────────────────────────────────────────
\section{Arquitetura da Ferramenta Demonstrativa}
\label{sec:arquitetura}

Para ilustrar a aplicação prática das diretrizes acima, foi construído o \textit{MikroTik Log Anonymizer}, disponibilizado via interface web em Streamlit. A ferramenta adota uma arquitetura em três estágios, resumida no \autoref{qdr:arquitetura}: detecção de IPs e MACs pelo Presidio Analyzer (recognizer nativo e customizado, respectivamente), detecção de nomes de usuário por expressão regular e transformação por meio dos operadores do Presidio AnonymizerEngine.

\begin{quadro}[htb]
    \caption{\label{qdr:arquitetura}Pipeline de processamento do MikroTik Log Anonymizer}
    \begin{tabular}{|p{1.5cm}|p{3.5cm}|p{8.8cm}|}
        \hline
        \textbf{Estágio} & \textbf{Componente} & \textbf{Responsabilidade} \\
        \hline
        1a & Presidio Analyzer & Detecção de endereços IP pelo recognizer nativo \texttt{IP\_ADDRESS} (score 0,92, reforçado por palavras-chave contextuais) e de endereços MAC pelo \textit{PatternRecognizer} customizado (score 0,95) \\
        \hline
        1b & Detecção por regex & Localização de nomes de usuário em eventos de autenticação MikroTik, categoria não coberta pelos recognizers padrão do Presidio \\
        \hline
        2 & Presidio Anonymizer & Aplicação dos operadores de transformação (hash, mask, redact, replace) sobre as entidades detectadas nos estágios anteriores \\
        \hline
        -- & Interface Streamlit & Upload do arquivo, configuração dos operadores e download do log anonimizado \\
        \hline
    \end{tabular}
    \legend{Fonte: Própria dos autores (2025).}
\end{quadro}

A combinação é deliberada: os recognizers do Presidio Analyzer oferecem detecção configurável e com pontuação de confiança para os identificadores de rede, enquanto o padrão textual dos eventos de autenticação do MikroTik -- específico da plataforma -- é mais eficientemente capturado por expressão regular dedicada. Os operadores do Presidio AnonymizerEngine, por sua vez, garantem a semântica e a configurabilidade das transformações aplicadas às entidades detectadas.

% ─────────────────────────────────────────────────────────────────────────────
\section{Motor de Anonimização (\texttt{core\_anonymizer.py})}
\label{sec:core}

O módulo \texttt{core\_anonymizer.py} concentra toda a lógica de desidentificação, organizada em estruturas de dados de configuração/resultado e padrões de detecção por expressão regular.

\begin{lstlisting}[language=Python, caption={Estruturas de dados de configuracao e resultado}, label={lst:dataclasses}]
@dataclass
class AnonymizerConfig:
    ip_mode: AnonymizationMode = "hash"   # hash | mask | redact | replace
    mac_mode: AnonymizationMode = "mask"
    redact_usernames: bool = True
    hash_type: str = "sha256"             # sha256 | sha512 | md5
    mask_char: str = "*"

@dataclass
class ProcessingResult:
    original_lines: int = 0
    processed_lines: int = 0
    ips_found: int = 0
    macs_found: int = 0
    usernames_found: int = 0
    elapsed_seconds: float = 0.0
    sample_before: list = field(default_factory=list)
    sample_after: list  = field(default_factory=list)
    output_bytes: bytes = b""
\end{lstlisting}

O tipo \texttt{AnonymizationMode} é um \texttt{Literal} Python que restringe os valores aceitos a \texttt{"hash"}, \texttt{"mask"}, \texttt{"redact"} e \texttt{"replace"}, garantindo validação estática por ferramentas como mypy. Como os endereços IP e MAC são detectados pelos recognizers do Presidio Analyzer (\autoref{sec:mac-recognizer}), a única expressão regular própria do módulo cobre os padrões de nome de usuário em eventos de autenticação MikroTik, categoria não coberta pelos recognizers padrão do framework:

\begin{lstlisting}[language=Python, caption={Padroes de deteccao de username por expressao regular}, label={lst:regex}]
# Padroes de username em eventos MikroTik
_USERNAME_PATTERNS = [
    (re.compile(r"\buser\s+(\w+)\s+logged\b", re.IGNORECASE),
     "user <USUARIO> logged"),
    (re.compile(r"\bchanged\s+by\s+(\w+)\b", re.IGNORECASE),
     "changed by <USUARIO>"),
    (re.compile(r"\bfailed\s+login\s+attempt\s+for\s+(\w+)\b",
     re.IGNORECASE), "failed login attempt for <USUARIO>"),
]
\end{lstlisting}

A classe \texttt{MikroTikAnonymizer} implementa o método \texttt{process\_file}, que itera sobre todas as linhas do log, aplicando sequencialmente os três estágios de anonimização:

\begin{lstlisting}[language=Python, caption={Pipeline de processamento linha a linha}, label={lst:pipeline}]
for i, line in enumerate(lines):
    # 1. Usernames -- Art. 6 III (Minimizacao)
    if config.redact_usernames:
        for pat, repl in _USERNAME_PATTERNS:
            new_line, n = pat.subn(repl, line)
            if n:
                line = new_line
                result.usernames_found += n

    # 2. IPs e MACs -- Art. 6 VII, VIII e IX
    #    (Seguranca / Prevencao / Nao discriminacao)
    analyzer_results = analyzer.analyze(
        text=line, entities=["IP_ADDRESS", "MAC_ADDRESS"], language="en"
    )
    if analyzer_results:
        result.ips_found += sum(
            1 for r in analyzer_results if r.entity_type == "IP_ADDRESS")
        result.macs_found += sum(
            1 for r in analyzer_results if r.entity_type == "MAC_ADDRESS")
        line = anonymizer.anonymize(
            text=line, analyzer_results=analyzer_results,
            operators=operators
        ).text

    out_lines.append(line)
\end{lstlisting}

A ordem dos estágios é deliberada: usernames são processados primeiro por expressão regular, antes de o restante da linha ser submetido ao Presidio Analyzer, evitando que fragmentos de nomes de usuário sejam erroneamente capturados pelos recognizers de IP ou MAC.

% ─────────────────────────────────────────────────────────────────────────────
\section{Recognizers do Presidio: IP Nativo e MAC Customizado}
\label{sec:mac-recognizer}

Para endereços IP, utiliza-se o recognizer nativo \texttt{IP\_ADDRESS} do Presidio Analyzer, com a lista de palavras-chave contextuais (\texttt{context}) ampliada para termos comuns em logs de rede -- como \texttt{"src"}, \texttt{"dst"}, \texttt{"ip"} e \texttt{"endereço"} --, elevando a pontuação de confiança das detecções para 0,92 nos testes realizados com o log real.

Já para endereços MAC, constatou-se que o Microsoft Presidio, em sua versão padrão (\texttt{presidio-analyzer >= 2.2.355}), não inclui um recognizer nativo -- resultado técnico relevante do projeto. Esse tipo de dado é tratado como PII de alta sensibilidade pela LGPD (art.\ 6°, IX), pois permite o perfilamento de dispositivos específicos ao longo do tempo. A implementação do recognizer customizado segue a interface \texttt{PatternRecognizer} do Presidio e é registrada no \texttt{AnalyzerEngine} antes do uso, ao lado do recognizer nativo de IP:

\begin{lstlisting}[language=Python, caption={Recognizer customizado para enderecos MAC}, label={lst:mac-recognizer}]
from presidio_analyzer import PatternRecognizer, Pattern

mac_pattern = Pattern(
    name="MAC_ADDRESS",
    regex=r"\b([0-9A-Fa-f]{2}:){5}[0-9A-Fa-f]{2}\b",
    score=0.95
)

mac_recognizer = PatternRecognizer(
    supported_entity="MAC_ADDRESS",
    patterns=[mac_pattern],
    context=["mac", "address", "device", "src-mac", "dst-mac"]
)
# Registrar no AnalyzerEngine antes do uso:
# analyzer.registry.add_recognizer(mac_recognizer)
\end{lstlisting}

O \textit{score} de 0,95 foi definido empiricamente após testes com o log real: a expressão regular é suficientemente estrita para o formato MAC que falsos positivos não foram observados, justificando um score próximo ao máximo. As palavras de contexto ampliam a precisão em casos onde o formato apareça em texto narrativo.

% ─────────────────────────────────────────────────────────────────────────────
\section{Operadores de Anonimização}
\label{sec:operadores}

Os operadores disponibilizados ao usuário na interface são configurados como objetos \texttt{OperatorConfig} do Presidio AnonymizerEngine, mapeados por tipo de entidade detectada. Para IPs, quatro modos são oferecidos:

\begin{lstlisting}[language=Python, caption={Fabrica de operadores Presidio para enderecos IP}, label={lst:op-ip}]
from presidio_anonymizer.entities import OperatorConfig

def _make_ip_operator(mode: str, hash_type: str) -> OperatorConfig:
    if mode == "hash":
        # Pseudonimizacao deterministica: mesmo IP -> mesmo hash
        return OperatorConfig("hash", {"hash_type": hash_type})
    if mode == "mask":
        return OperatorConfig("mask", {
            "masking_char": "*", "chars_to_mask": 9, "from_end": False
        })
    if mode == "redact":
        return OperatorConfig("redact", {})
    return OperatorConfig("replace", {"new_value": "<IP_ANONIMIZADO>"})
\end{lstlisting}

A propriedade determinística do operador \texttt{hash} é essencial para a análise de incidentes: o analista pode agrupar todos os eventos originados de um mesmo IP sem conhecer o endereço real. Para MACs, a mesma lógica de quatro modos é aplicada, preservando por padrão os três octetos finais do dispositivo:

\begin{lstlisting}[language=Python, caption={Fabrica de operadores Presidio para enderecos MAC}, label={lst:op-mac}]
def _make_mac_operator(mode: str) -> OperatorConfig:
    if mode == "hash":
        return OperatorConfig("hash", {"hash_type": "sha256"})
    if mode == "mask":
        # Mascara os 3 primeiros octetos (OUI/fabricante)
        # Preserva os 3 ultimos (identificador do dispositivo)
        return OperatorConfig("mask", {
            "masking_char": "*", "chars_to_mask": 9, "from_end": False
        })
    if mode == "redact":
        return OperatorConfig("redact", {})
    return OperatorConfig("replace", {"new_value": "<MAC_ANONIMIZADO>"})

operators = {
    "IP_ADDRESS": _make_ip_operator(config.ip_mode, config.hash_type),
    "MAC_ADDRESS": _make_mac_operator(config.mac_mode),
}
\end{lstlisting}

Os três primeiros octetos do MAC formam o OUI (\textit{Organizationally Unique Identifier}), que identifica o fabricante e por si só não identifica o titular; os três finais identificam o dispositivo específico e constituem o dado pessoal sensível. O resultado \texttt{***:***:***:XX:XX:XX} remove a possibilidade de perfilamento sem eliminar a utilidade do dado para análise de topologia de rede.

% ─────────────────────────────────────────────────────────────────────────────
\section{Pseudonimização por Hash com \textit{Salt}}
\label{sec:hash}

Para cenários que exigem garantia adicional contra ataques de dicionário ou pré-computação (\textit{rainbow tables}), recomenda-se complementar o hash com um \textit{salt} -- valor aleatório concatenado ao dado antes do resumo criptográfico -- mantido separadamente dos logs, em cofre de segredos com acesso restrito \cite{dwork2014differential}.

\begin{lstlisting}[language=Python, caption={Pseudonimizacao de IP com SHA-256 e salt}, label={lst:hash-ip}]
import hashlib, secrets

SALT = secrets.token_hex(32)  # armazenado separadamente dos logs

def pseudonimizar_ip(ip: str) -> str:
    """Mesmo IP + salt sempre gera o mesmo hash (correlacao
    preservada); sem o salt, o hash e irreversivel por dicionario."""
    return hashlib.sha256((SALT + ip).encode("utf-8")).hexdigest()
\end{lstlisting}

% ─────────────────────────────────────────────────────────────────────────────
\section{Generalização Temporal}
\label{sec:generalizacao-temporal}

A generalização temporal trunca a precisão dos \textit{timestamps}, reduzindo o risco de quasi-identificação por combinação de horário com outros atributos. Para compartilhamento externo, recomenda-se truncar ao nível de hora ou dia.

\begin{lstlisting}[language=Python, caption={Generalizacao temporal de timestamps em logs MikroTik}, label={lst:timestamp}]
_TS_RE = re.compile(r"([a-z]{3}/\d{2}\s+)(\d{2}:\d{2}:\d{2})", re.IGNORECASE)

def generalizar_timestamp(linha: str, nivel: str = "hora") -> str:
    """nivel='hora' -> preserva hh:00:00; nivel='dia' -> remove horario."""
    def substituir(m):
        if nivel == "hora":
            return m.group(1) + m.group(2).split(":")[0] + ":00:00"
        return m.group(1).strip()
    return _TS_RE.sub(substituir, linha)
\end{lstlisting}

% ─────────────────────────────────────────────────────────────────────────────
\section{Interface Web (\texttt{app.py})}
\label{sec:interface}

A interface web foi desenvolvida em Streamlit e estruturada em quatro áreas funcionais:

\begin{alineas}
    \item \textbf{Painel lateral de configuração}: o usuário seleciona, de forma independente, o método de anonimização para IPs (com seleção do algoritmo de hash quando aplicável), para MACs e o comportamento para usernames. O painel exibe também os artigos LGPD cobertos por cada configuração, servindo como guia de conformidade embutido na interface;
    \item \textbf{Área de upload}: aceita arquivos \texttt{.log}, \texttt{.txt} e \texttt{.zip}. Arquivos comprimidos são descompactados em memória (RAM), sem gravação em disco, antes do processamento;
    \item \textbf{Painel de resultados}: exibe métricas da execução (linhas processadas, contagens por tipo de PII, tempo decorrido e tamanho do log anonimizado) e um preview lado a lado de até 10 linhas com PII detectada, antes e depois da anonimização, permitindo auditoria visual imediata;
    \item \textbf{Download}: o log anonimizado é disponibilizado para download direto no navegador, com nomenclatura automática (\texttt{<nome-original>\_lgpd\_anonimizado.log}).
\end{alineas}

O processamento assíncrono com barra de progresso (atualizada a cada 0,5\% do arquivo) mantém a interface responsiva durante o processamento de arquivos grandes. O motor é carregado com cache (\texttt{@st.cache\_resource}), evitando reinicializações desnecessárias entre execuções consecutivas.

% ─────────────────────────────────────────────────────────────────────────────
\section{Resultados da Validação}
\label{sec:presidio-pipeline}

O pipeline completo foi validado com um log real de roteador MikroTik em produção, contendo 738.846 linhas e cerca de 50~MB de dados brutos. O processamento anonimizou 258.688 IPs, 148.120 MACs e 73.867 usernames em aproximadamente 9 segundos em hardware convencional, demonstrando viabilidade operacional para uso cotidiano. A decisão de preservar os \textit{timestamps} originais mereceu análise específica: embora possam contribuir para reidentificação em combinação com outros atributos (\autoref{sec:reidentificacao}), sua supressão eliminaria a utilidade primária dos logs para análise de incidentes -- optou-se por mantê-los para uso interno, com generalização temporal recomendada para cenários de compartilhamento externo.

Como limitação, a ferramenta não trata o risco de reidentificação por combinação de quasi-identificadores (horário + tipo de evento + segmento de rede). Para cenários de alta sensibilidade, recomenda-se complementar com técnicas de privacidade diferencial, implementáveis via \textit{SmartNoise} ou \textit{Google Differential Privacy Library} \cite{dwork2014differential}.

% ─────────────────────────────────────────────────────────────────────────────
\section{Diretrizes para Aplicação}
\label{sec:diretrizes-aplicacao}

Com base nas técnicas descritas, propõem-se as seguintes diretrizes para a organização que adotar as políticas deste projeto:

\begin{alineas}
    \item Aplicar anonimização automaticamente no momento da coleta, sempre que a finalidade primária do log não exigir a identidade do titular -- eliminando o risco na origem;
    \item Para logs que precisam preservar a capacidade de correlação (uso interno de segurança), adotar pseudonimização por hash com \textit{salt}; para compartilhamento externo, aplicar supressão ou generalização;
    \item Manter separação física ou lógica entre o \textit{salt} e os logs pseudonimizados -- o acesso ao \textit{salt} deve ser restrito ao administrador de segurança;
    \item Documentar o processo de anonimização aplicado a cada categoria de log, incluindo o algoritmo utilizado, os campos tratados e a data de aplicação, para fins de demonstração de conformidade (art.\ 6°, X, LGPD -- responsabilização), integrando o registro de atividades de tratamento do \autoref{sec:documentacao};
    \item Revisar anualmente a efetividade das técnicas adotadas, considerando a evolução das capacidades técnicas de reidentificação.
\end{alineas}